\begin{problem}[第34届物理决赛][荡秋千能量分析]{82}
    熟练的荡秋千的人能够通过在秋千板上适时站起和蹲下使秋千越荡越高。一质量为 $m$ 的人荡一架底板和摆杆均为刚性的秋千，底板和摆杆的质量均可忽略，假定人的质量集中在其质心。人在秋千上每次完全站起时其质心距悬点 $O$ 的距离为 $l$，完全蹲下时此距离变为 $l+d$。实际上，人在秋千上站起和蹲下过程都是在一段时间内完成的。作为一个简单的模型，假设人在第1个最高点 $A$ 点从完全站立的姿势迅速完全下蹲，然后荡至最低点 $B$，$A$ 与 $B$ 的高度差为 $h_{1}$；随后他在 $B$ 点迅速完全站起（且最终径向速度为零），继而随秋千荡至第2个最高点 $C$，这一过程中该人质心运动的轨迹如图所示。此后人以同样的方式回荡，重复前述过程，荡向第3、4等最高点。假设人在站起和蹲下的过程中，人与秋千的相互作用力始终与摆杆平行。以最低点 $B$ 为重力势能零点。
    \begin{subproblem}
        \begin{subsubproblem}
            假定在始终完全蹲下和始终完全站立过程中没有机械能损失，求该人质心在 $A\rightarrow A'\rightarrow B\rightarrow B'\rightarrow C$ 各个阶段的机械能及其变化；
        \end{subsubproblem}
        \begin{subsubproblem}
            假定在始终完全蹲下和始终完全站立过程中的机械能损失 $\Delta E$ 与过程前后高度差的绝对值 $\Delta h$ 的关系分别为
            $\Delta E = k_{1} m g (h_{0} + \Delta h)$, $0<k_{1}<1$, 始终完全蹲下； $\Delta E = k_{2} m g (h_{0}' + \Delta h)$, $0<k_{2}<1$, 始终完全站立。
            这里，$k_{1}$、$k_{2}$、$h_{0}$ 和 $h_{0}'$ 是常量，$g$ 是重力加速度的大小。求
            \begin{subsubsubproblem}
                \begin{subsubsubsubproblem}
                    相对于 $B$ 点，第 $n$ 个最高点的高度 $h_{n}$ 与第 $n+1$ 个最高点的高度 $h_{n+1}$ 之间的关系；
                \end{subsubsubsubproblem}
                \begin{subsubsubsubproblem}
                    $h_{n}$ 与 $h_{1}$ 之间的关系式和 $h_{n+1} - h_{n}$ 与 $h_{1}$ 之间的关系式。
                \end{subsubsubsubproblem}
            \end{subsubsubproblem}
        \end{subsubproblem}
    \end{subproblem}
\end{problem}